\chapter{The matrices of gsm}
\label{ch:gsm}

Gnascor reads two answers side by side as a branch pair \cite{src:gnascor}. The pair is a state of a finite
machine, and each transition of that machine carries a mnemonic. The table \texttt{gnascor\_asm\_lng.tsv} holds 86
rows: 4 of the branch pair, 63 of the state-transition matrix, 12 of the high-energy matrix, 2 of the shifts and 5
of the physics and error mnemonics. They name 28 mnemonics.

\section{States and transitions}

\begin{definition}
A state with a pair is one of lead, rite, dual and void, with the bits $(1,0)$, $(0,1)$, $(1,1)$ and $(0,0)$. Every
other state holds no pair bit and writes $-$ for each side. A transition is a row: the state it leaves, the state
it reaches, the bits of each, an operation for each side, an arrow and a mnemonic.
\end{definition}

\begin{definition}
The operation of a side, from its bit $p$ before and $q$ after, is $-$ where $p$ or $q$ is $-$, $=$ where $p = q$,
and $q$ where $p \neq q$. The arrow is \texttt{*>} where a side is $-$, \texttt{x>} where both sides change, and
\texttt{->} where a side stays.
\end{definition}

A branch pair row has two states for its sides and no transition. Each of its sides is a state of the pair below,
and each state is its two bits. A pair of branches at once is written \texttt{[ ]}, and a state over two cycles
is written \texttt{< >}.

\section{The laws}

Each law below is a conjunction over the 86 rows, and the kernel decides each one.

\begin{theorem}[\texttt{side\_law}, \texttt{arrow\_law}]
In every row that is a transition, each side's operation follows from its two bits, and the arrow follows from the
two operations.
\end{theorem}

\begin{theorem}[\texttt{state\_law}]
In every row, the branch pair included, the name of each state fixes its bits.
\end{theorem}

\begin{theorem}[\texttt{mnemonic\_law}]
One table gives one transition at most one mnemonic.
\end{theorem}

\section{The transpiler form of a mnemonic}

A transition that gives each side a bit does one of two things to that side: it keeps the bit or it inverts it.
The two are ladder anchors, copy and invert. Let $\kappa(=) = \mathrm{copy}$ and $\kappa(0) = \kappa(1) =
\mathrm{invert}$.

\begin{theorem}[\texttt{form\_law}]
In every transition, each side with a pair bit reaches the bit $R_{\kappa(o)}(p)$, where $o$ is its operation and
$p$ its bit before.
\end{theorem}

The transpiler form of a transition is the bridge word of each side's anchor: \texttt{set} for copy and
\texttt{not} for invert. \texttt{lng\_asm\_form} of \texttt{lng.c} derives it, and the language check reports each
mnemonic with no transition that gives each side a bit. 13 of the 28 mnemonics have a form. The other 15 name
only transitions where a side holds no pair bit, and they remain holes.
